Large language models have made synthetic data inexpensive, but they can still be biased. When real data are scarce, researchers must decide how much of the real sample to use for direct estimation and how much to spend calibrating the generator. We study this allocation problem and derive the marginal-value condition for the optimal calibration size: $(v_n + c x^{-2\beta})^2 = 2\beta a c x^{-(2\beta+1)}$, where $x$ is the calibration size, $v_n$ is synthetic sampling variance, $a$ is the real-estimator variance constant, and $c x^{-2\beta}$ is squared synthetic bias. This condition has no universal fixed-share solution and yields five allocation regimes. In the main regime, $m_n \asymp n$ and $\beta > 1/2$, the optimal calibration size grows only as $n^{2/(2\beta+1)}$, so the calibration share vanishes; rules that keep a positive fixed fraction of real observations in calibration over-calibrate asymptotically. We also propose an adaptive grid estimator with an oracle inequality and a safety check, $xB_n(x)<a$, that falls back to all-real estimation when synthetic data cannot help. Across approximately $10^7$ Monte Carlo replications, the corrected oracle reduces MSE by about $47\%$ relative to all-real estimation in the headline cell, compared with about $22\%$ for a fixed-share baseline. For inference, naive Wald intervals around the MSE-optimal estimator can undercover, while validation debiasing substantially restores coverage. These results provide guidance for when synthetic data are useful, how much real data should be used to make them reliable, and how to conduct valid inference with synthetic-data-augmented estimators.
